% Kato，物理页139--168；原书134--163。
\providecommand{\kaOG}{\mathcal O(G)}
\providecommand{\kaOm}{\Omega^1(G)}
\providecommand{\kaco}{\operatorname{coLie}(G)}
\providecommand{\kaJ}{J_{\infty,\bar K/K}}
\providecommand{\kaB}{B_{\infty,\bar K/K}}
\providecommand{\kaTr}{\operatorname{Trace}}
\origpage{139}{134}
由下面的引理，定理 2.1.7 可归结为定理 2.2.7。
\begin{statement}{引理 2.2.8}
对任意 $n\geq0$，下图交换\translatornote{13}：
\[
\resizebox{\linewidth}{!}{$\begin{tikzcd}[ampersand replacement=\&,column sep=large,row sep=large]
\kaOG^\times\arrow[r,"{\xi_n}"]\arrow[d,"{\mathrm{dlog}}"']\&
\mathcal O_{K_n}^\times\arrow[r]\&H^1(K_n,\mathbb Z_p(1))\arrow[r]\arrow[d]\&
H^1(K_n,T^{\otimes(-r)}(1)/\pi^n)\arrow[d]\\
\kaOm\arrow[rr,"{\mu_{\xi,n}}"']\&\&H^1(K_n,\kaJ)\arrow[r]\&
H^1(K_n,(T^{\otimes(-r)}\otimes_{\mathcal O_F}\kaJ)/\pi^n).
\end{tikzcd}$}
\]
\end{statement}
\noindent\textbf{2.2.9.}
证明 2.2.8 之前，先明确描述映射 $\mu_{\xi,n}$。考虑由
$\kaOG\xrightarrow{\xi_n}\mathcal O_{K_n}\longrightarrow\widehat{\mathcal O}_{\bar K}$
得到的映射
\[
\kaOG\otimes_{\mathcal O_K}\kaB\longrightarrow\widehat{\mathcal O}_{\bar K}.
\]
设 $\bar S_n$ 为满足下述性质的 $p$ 进完备且分离的拓扑环：对任意 $i\geq1$，
$\Spec(\bar S_n/p^i\bar S_n)$ 是
$\Spec(\mathcal O_{\bar K}/p^i\mathcal O_{\bar K})$ 在
$\Spec((\kaOG\otimes_{\mathcal O_K}\kaB)/p^i)$ 中的
\term{PD包络}{PD envelope}{enveloppe à puissances divisées}。
于是 $\bar S_n$ 带有 $\Gal(\bar K/K_n)$ 的自然作用，并具有如下结构：存在一个保持
\term{分次幂结构}{divided power structure}{structure à puissances divisées}的 $\kaB$ 上同构
\[
\varprojlim_m B_m\langle t\rangle\xrightarrow{\sim}\bar S_n,
\qquad t\longmapsto h\otimes1-1\otimes x.
\]
这里 $\langle\ \rangle$ 表示 PD 多项式环，$t$ 是不定元；我们固定了
\term{增广理想}{augmentation ideal}{idéal d'augmentation}
$\Ker(\kaOG\to\mathcal O_K)$ 的一个生成元 $h$，以及 $\kaB$ 中的一个元素 $x$，使其在
$\widehat{\mathcal O}_{\bar K}$ 中的像等于 $h(\xi_n)$。令
$J_{\bar S_n}=\Ker(\bar S_n\to\mathcal O_{\bar K})$。则序列
\[
0\longrightarrow\kaJ\longrightarrow J_{\bar S_n}
\xrightarrow{d}\bar S_n\otimes_{\kaOG}\kaOm\longrightarrow0
\]
正合。这由 $d(h\otimes1-1\otimes x)=1\otimes dh$，以及 $dh$ 是可逆
$\kaOG$ 模 $\kaOm$ 的一组基这一事实得到。

\origpage{140}{135}
映射 $\mu_{\xi,n}$ 就是这一序列在上同调中给出的\term{连接映射}{connecting map}{homomorphisme de connexion}。

\noindent\textbf{2.2.10.} 现在证明 2.2.8。右边方块的交换性显然，因此只考虑左边的方块。
设 $h\in\kaOG^\times$。任取 $m\geq0$，并取
$v\in\mathcal O_{\bar K}^\times$ 使 $v^{p^m}=h(\xi_n)$。则 $h$ 在
$H^1(K_n,\mathbb Z/p^m(1))$ 中的像由余循环
\[
\Gal(\bar K/K_n)\longrightarrow\mathbb Z/p^m(1),
\qquad\sigma\longmapsto\frac{\sigma(v)}v
\]
表示。取 $\tilde v\in\kaB^\times$，使其在
$\widehat{\mathcal O}_{\bar K}^\times$ 中的像等于 $v$。回顾映射
\[
\mathbb Z/p^m(1)\longrightarrow\kaJ/p^m,
\qquad\alpha\longmapsto\log(\tilde\alpha^{p^m})
\]
的定义，其中 $\tilde\alpha$ 是 $1$ 的 $p^m$ 次单位根在
$(\kaB/p^m)^\times$ 中的一个提升（1.1.5），可知余循环
$\sigma\mapsto\sigma(v)/v$ 在 $H^1(K_n,\kaJ/p^m)$ 中的像由下面的余循环表示：
\begin{align*}
\sigma&\longmapsto\log\left(\left(\frac{\sigma(\tilde v)}{\tilde v}\right)^{p^m}\right)\\
&=\log\left(\frac{\sigma(\tilde v^{p^m}h^{-1})}{\tilde v^{p^m}h^{-1}}\right)
=\sigma\bigl(\log(\tilde v^{p^m}h^{-1})\bigr)-\log(\tilde v^{p^m}h^{-1}).
\end{align*}
这里 $\log(\tilde v^{p^m}h^{-1})\in J_{\bar S_n}$ 有定义，因为
$\tilde v^{p^m}h^{-1}\in\bar S_n^\times$ 在
$\widehat{\mathcal O}_{\bar K}^\times$ 中的像为 $1$。由
\[
d\log(\tilde v^{p^m}h^{-1})\equiv-\frac{dh}{h}\pmod{p^m}
\]
及 2.2.9，最后这个余循环的类等于 $\mu_{\xi,n}(dh/h)$ 在
$H^1(K_n,\kaJ/p^m)$ 中的像。\translatornoteref{13}

\section*{§2.3. 定理 2.2.7 的证明}
\addcontentsline{toc}{section}{§2.3. 定理 2.2.7 的证明}
下面证明中 $r=1$ 的情形由栗原将人给出。

\noindent\textbf{2.3.1.} 以下把 $\kaco^{\otimes r}$ 看作嵌入在
\[
H^0\bigl(K,T^{\otimes(-r)}\otimes_{\mathcal O_F}
(\kaJ^{[r]}/\kaJ^{[r+1]})\bigr)
\]
中（2.2.4）。固定 $n\geq1$。对于满足
$(r-1)!^{-1}\pi^{-nr}c\in\mathcal O_K$ 的整数 $c$，以及 $m\geq n$，定义映射
\origpage{141}{136}
\[
\alpha_{c,m},\ \beta_{c,m},\ \tau_{c,m}:
S_m\otimes_{\kaOG}\kaOm\longrightarrow
H^1\bigl(K_n,(T^{\otimes(-r)}\otimes_{\mathcal O_F}
\kaJ/\kaJ^{[r+1]})/\pi^m\bigr)
\]
如下。令 $\alpha_{c,m}$ 为以下复合映射的 $c$ 倍：
\begin{gather*}
S_m\otimes_{\kaOG}\kaOm\xrightarrow{\mu_{r,\xi,m}}
H^1\bigl(K_m,(T^{\otimes(-r)}\otimes_{\mathcal O_F}
\kaJ/\kaJ^{[r+1]})/\pi^m\bigr)\\
\xrightarrow{\mathrm{trace}}
H^1\bigl(K_n,(T^{\otimes(-r)}\otimes_{\mathcal O_F}
\kaJ/\kaJ^{[r+1]})/\pi^m\bigr).
\end{gather*}
令 $\beta_{c,m}$ 由下式定义：
\[
\beta_{c,m}(\eta)=
\bigl((r-1)!^{-1}\pi^{-nr}c\bigr)\omega^{\otimes r}
\otimes\left\{\left(\frac d\omega\right)^{r-1}
\left(\frac1\omega\kaTr_{m,n}(\eta)\right)\right\}(\xi_n)
\cup\log(\chi_G),
\]
其中 $\omega$ 是 $\kaco$ 的任意一组 $\mathcal O_K$ 基；
$\kaTr_{m,n}:S_m\otimes_{\kaOG}\kaOm\to S_n\otimes_{\kaOG}\kaOm$
是下文定义的迹映射；而 $\log(\chi_G)$ 是
\[
H^1(K,\mathcal O_F)=\Hom_{\mathrm{cont}}(\Gal(\bar K/K),\mathcal O_F)
\]
中的元素，定义为复合映射
\[
\Gal(\bar K/K)\xrightarrow{\chi_G}\mathcal O_F^\times
\xrightarrow{\log}\mathcal O_F,
\]
其中 $\chi_G$ 是在 $T$ 上的作用。\translatornote{14}令 $\tau_{c,m}=\alpha_{c,m}+\beta_{c,m}$。

现给出上述 $\kaTr_{m,n}$ 的定义，并稍作补充。有环的交换图
\[
\begin{tikzcd}[ampersand replacement=\&,column sep=large,row sep=large]
\kaOG\arrow[r]\arrow[d,"{[\pi^{m-n}]^*}"']\&S_n\arrow[r,"{\xi_n}"]\arrow[d,"{[\pi^{m-n}]^*}"]\&\mathcal O_{K_n}\arrow[d]\\
\kaOG\arrow[r]\&S_m\arrow[r,"{\xi_m}"']\&\mathcal O_{K_m}.
\end{tikzcd}
\]
图中各方块为\term{推出}{push-out}{somme amalgamée}，其中
$[\pi^{m-n}]^*$ 是由 $[\pi^{m-n}]:G\to G$ 诱导的同态。
与这些竖直箭头相伴的迹映射均记为 $\kaTr_{m,n}$，并有交换图
\origpage{142}{137}
\[
\begin{tikzcd}[ampersand replacement=\&,column sep=large,row sep=large]
\kaOG\arrow[r]\arrow[d,"{\kaTr_{m,n}}"']\&S_m\arrow[r,"{\xi_m}"]\arrow[d,"{\kaTr_{m,n}}"]\&\mathcal O_{K_m}\arrow[d,"{\kaTr_{m,n}}"]\\
\kaOG\arrow[r]\&S_n\arrow[r,"{\xi_n}"']\&\mathcal O_{K_n}.
\end{tikzcd}
\]
容易看出 $\kaTr_{m,n}(\kaOG)\subset\pi^{m-n}\kaOG$，因而
$\kaTr_{m,n}(S_m)\subset\pi^{m-n}S_n$。所需映射
$\kaTr_{m,n}:S_m\otimes_{\kaOG}\kaOm\to S_n\otimes_{\kaOG}\kaOm$
定义为
\[
\kaTr_{m,n}(x\otimes\omega)=
\bigl(\pi^{n-m}\kaTr_{m,n}(x)\bigr)\otimes\omega
\quad(x\in S_m,\ \omega\in\kaco).
\]
注意 $\pi^{n-m}\kaTr_{m,n}(x)\in S_n$。这个定义是自然的，因为有投影公式
\[
\kaTr_{m,n}(x\otimes[\pi^{m-n}]^*\omega)=\kaTr_{m,n}(x)\otimes\omega
\quad(x\in S_m,\ \omega\in\kaco),
\]
其中 $[\pi^{m-n}]^*\omega=\pi^{m-n}\omega$。

现在，定理 2.2.7 可由下面的命题 2.3.2 及如下事实推出：
$\log(\chi_G)$ 与 $\log(\chi_{\mathrm{cyclo}})$ 在 $H^1(K,\mathbb C_p)$ 中的像相同。
后一个事实来自 $V\otimes_F\mathbb C_p\simeq\mathbb C_p(1)$。
\begin{statement}{命题 2.3.2}
存在整数 $c$，满足 $(r-1)!^{-1}\pi^{-nr}c\in\mathcal O_K$，使得对所有 $m\geq n$
都有 $\tau_{c,m}=0$。这里的 $c$ 可取为与 $m$ 无关。
\end{statement}
容易看出 $\alpha_{c,m}$ 与 $\beta_{c,m}$，因而 $\tau_{c,m}$，在
$J_{S_m}^{[r]}\otimes_{\kaOG}\kaOm$ 上为零。
下面的 2.3.3 是证明 2.3.2 的关键。
\begin{statement}{引理 2.3.3}
存在非零整数 $c$，满足 $(r-1)!^{-1}\pi^{-nr}c\in\mathcal O_K$，并且对所有 $m\geq n$，
$\tau_{c,m}$ 都消去 $\ell_{\omega,\xi,m}^{r-1}S_m\otimes_{\kaOG}\kaOm$。
\end{statement}
这个引理将由下述引理推出（见 2.3.8）。
\begin{statement}{引理 2.3.4}
映射
\[
S_m\otimes_{\kaOG}\kaOm\xrightarrow{\mu_{r,\xi,m}}
H^1\bigl(K_m,(T^{\otimes(-r)}\otimes_{\mathcal O_F}
\kaJ/\kaJ^{[r+1]})/\pi^m\bigr)
\]
\origpage{143}{138}
将 $\ell_{\omega,\xi,m}^{r-1}h\otimes\omega$
（$h\in S_m$，$\omega\in\kaco$）映为
\[
-\omega^{\otimes r}\otimes h(\xi_m)\cup\chi_m,
\]
其中 $\chi_m$ 是
$H^1(K_m,\mathcal O_F/\pi^m)=\Hom_{\mathrm{cont}}(\Gal(\bar K/K_m),\mathcal O_F/\pi^m)$
中由下式定义的元素：
\[
\chi_m(\sigma)=\pi^{-m}(\chi_G(\sigma)-1)\bmod\pi^m
\quad(\sigma\in\Gal(\bar K/K_m)).
\]
\end{statement}
我们将利用下面的引理 2.3.6 来证明 2.3.4。
\begin{statement}{引理 2.3.5}
正合序列
\[
0\longrightarrow\kaJ\xrightarrow{\pi^m}\kaJ
\longrightarrow\kaJ/\pi^m\longrightarrow0
\]
的连接映射将 $\ell_{\omega,\xi,m}\in H^0(K_m,\kaJ/\pi^m)$ 映为
$-\mu_{\xi,m}(\omega)\in H^1(K_m,\kaJ)$。\footnote{译注：原书此处印作 $K_n$；依本段定义及下文证明改为 $K_m$。}
\end{statement}
\begin{proof}
取 $x\in J_{\bar S_m}$ 使 $dx=\omega$。则 $\mu_{\xi,m}(\omega)\in H^1(K_m,\kaJ)$
由余循环
\[
\Gal(\bar K/K_m)\longrightarrow\kaJ,
\qquad\sigma\longmapsto\sigma(x)-x
\]
表示（2.2.9）。令 $y=\ell_{\omega,\xi,m}-\pi^m x$。则 $dy=0$，故 $y\in\kaJ$，且
$y\bmod\pi^m=\ell_{\omega,\xi,m}\bmod\pi^m$。于是余循环
\[
\Gal(\bar K/K_m)\longrightarrow\kaJ,
\qquad\sigma\longmapsto\pi^{-m}(\sigma(y)-y)
\]
表示 $\ell_{\omega,\xi,m}\in H^0(K_m,\kaJ/\pi^m)$ 在连接映射下的像。
然而 $\pi^{-m}(\sigma(y)-y)=x-\sigma(x)$。
\end{proof}
\begin{statement}{引理 2.3.6}
在 $H^1(K_m,\kaJ/\pi^m)$ 中有
\[
\ell_{\omega,\xi,m}\cup\chi_m=-\mu_{\xi,m}(\omega).
\]
这里把 $\ell_{\omega,\xi,m}$ 看作 $H^0(K_m,\kaJ/\pi^m)$ 中的元素。
\end{statement}
\begin{proof}
令 $X=\kaco\otimes_{\mathcal O_F}T$，$Y=\kaJ$。考虑下列正合序列的交换图：
\origpage{144}{139}
\[
\begin{tikzcd}[ampersand replacement=\&,column sep=large,row sep=large]
0\arrow[r]\&X/\pi^m\arrow[r,"{\pi^m}"]\arrow[d]\&X/\pi^{2m}\arrow[r]\arrow[d]\&X/\pi^m\arrow[r]\arrow[d]\&0\\
0\arrow[r]\&Y/\pi^m\arrow[r,"{\pi^m}"']\&Y/\pi^{2m}\arrow[r]\&Y/\pi^m\arrow[r]\&0.
\end{tikzcd}
\]
其中竖直箭头是配对 $\omega\otimes\xi\mapsto\ell_{\omega,\xi,m}$。由此得到交换图
\[
\begin{tikzcd}[ampersand replacement=\&,column sep=large,row sep=large]
H^0(K_m,X/\pi^m)\arrow[r]\arrow[d]\&H^1(K_m,X/\pi^m)\arrow[d]\\
H^0(K_m,Y/\pi^m)\arrow[r]\&H^1(K_m,Y/\pi^m).
\end{tikzcd}
\]
容易看出，上方横箭头将
\[
\omega\otimes\xi_m\in\kaco\otimes H^0(K_m,T/\pi^m)=H^0(K_m,X/\pi^m)
\]
映为
$\omega\otimes\chi_m\cup\xi_m\in\kaco\otimes H^1(K_m,T/\pi^m)=H^1(K_m,X/\pi^m)$，
最后这个元素经右侧竖箭头映为 $\ell_{\omega,\xi,m}\cup\chi_m$。
另一方面，由 2.3.5，下方横箭头将 $\ell_{\omega,\xi,m}$ 映为
$-\mu_{\xi,m}(\omega)$。这就证明了 2.3.6。
\end{proof}
\noindent\textbf{2.3.7.} 利用 2.3.6 证明 2.3.4。在
$H^1(K_m,(T^{\otimes(-r)}\otimes_{\mathcal O_F}\kaJ/\kaJ^{[r+1]})/\pi^m)$ 中，
对于 $h\in S_m$ 和 $\kaco$ 的一组基 $\omega$，有
\begin{align*}
\mu_{r,\xi,m}(\ell_{\omega,\xi,m}^{r-1}h\otimes\omega)
&=(\xi^{\otimes r})^{-1}\otimes
\ell_{\omega,\xi,m}^{r-1}h(\xi_m)\mu_{\xi,m}(\omega)
&&\text{由 $\mu_{r,\xi,m}$ 的定义}\\
&=-\omega^{\otimes r}\otimes h(\xi_m)\chi_m
&&\text{由 2.3.6。}
\end{align*}
\noindent\textbf{2.3.8.} 由 2.3.4 推出 2.3.3。先考虑 $\alpha_{c,m}$。由于
\[
\log(\chi_G(\sigma))=(\chi_G(\sigma)-1)+\frac12(\chi_G(\sigma)-1)^2+\cdots,
\]
\footnote{译注：对数展开二次项的正号原书如此，与通常展开不符；此处照录。}
存在一个与 $m$ 无关的非零整数 $c$，使得
\begin{equation*}\tag{*}
\begin{gathered}
c\pi^{-m}\log(\chi_G(\sigma))\in\mathcal O_F,\\
c\pi^{-m}\log(\chi_G(\sigma))\equiv
c\pi^{-m}(\chi_G(\sigma)-1)\pmod{\pi^m}
\end{gathered}
\end{equation*}
对一切 $m\geq n$ 及 $\sigma\in\Gal(\bar K/K_m)$ 成立。
\origpage{145}{140}
对满足（*）的非零整数 $c$ 及 $h\in S_m$，由 2.3.4 有
\[
\alpha_{c,m}(\ell_{\omega,\xi,m}^{r-1}h\otimes\omega)
=-\omega^{\otimes r}\otimes\kaTr_{m,n}
\bigl(h(\xi_m)(c\pi^{-m}\log(\chi_G))\bigr).
\]
这里 $\kaTr_{m,n}$ 是迹映射
$H^1(K_m,\mathcal O_{\bar K}/\pi^m)\to H^1(K_n,\mathcal O_{\bar K}/\pi^m)$。
最后这个表达式的 $\pi$ 倍，等于
\[
-\omega^{\otimes r}\otimes
\bigl(\pi^{n-m}\kaTr_{m,n}(h(\xi_m))\bigr)(c\pi^{-n}\log(\chi_G))
\]
的 $\pi$ 倍；此处 $\kaTr_{m,n}$ 是迹映射
$\mathcal O_{K_m}\to\mathcal O_{K_n}$，其像包含在 $\pi^{m-n}\mathcal O_{K_n}$ 中。
理由是：对 $n<i\leq m$，迹映射
$H^1(K_i,\mathcal O_{\bar K})\to H^1(K_{i-1},\mathcal O_{\bar K})$
将
\[
\bigl(\pi^{i-m}\kaTr_{m,i}(h(\xi_m))\bigr)(c\pi^{1-i}\log(\chi_G))
\]
映为
\[
\bigl(\pi^{i-m}\kaTr_{m,i-1}(h(\xi_m))\bigr)(c\pi^{1-i}\log(\chi_G)),
\]
因为 $c\pi^{1-i}\log(\chi_G)$ 来自 $H^1(K_{i-1},\mathcal O_F)$。

再考虑 $\beta_{c,m}$。对 $h\in S_m$ 有
\[
\kaTr_{m,n}(\ell_{\omega,\xi,m}^{r-1}h\otimes\omega)
=\ell_{\omega,\xi,n}^{r-1}
\bigl(\pi^{n-m}\kaTr_{m,n}(h)\bigr)\otimes\omega,
\]
右侧的 $\kaTr_{m,n}$ 是与 $[\pi^{m-n}]^*:S_n\to S_m$ 相伴的迹映射（2.3.1）。
由 $d\ell_{\omega,\xi,m}=\pi^n\omega$\footnote{译注：此式左侧下标为 $m$、右侧指数为 $n$，原书如此；结合随后公式，疑有下标误植，此处照录。}，对任意 $g\in S_n$ 有
\[
\left(\frac d\omega\right)^{r-1}
(\ell_{\omega,\xi,n}^{r-1}g)
\equiv(r-1)!\pi^{n(r-1)}g\pmod{J_{S_n}}.
\]
由此得
\begin{align*}
&\left\{\left(\frac d\omega\right)^{r-1}
\left(\frac1\omega\kaTr_{m,n}(\ell_{\omega,\xi,m}^{r-1}h\otimes\omega)\right)\right\}(\xi_n)\\
&\hspace{2em}=(r-1)!\pi^{n(r-1)}\pi^{n-m}\kaTr_{m,n}(h(\xi_m)).
\end{align*}
因而由 2.3.4，对某个满足 $(r-1)!^{-1}\pi^{-nr}c\in\mathcal O_K$ 的非零整数 $c$，
$-\alpha_{c,m}(\ell_{\omega,\xi,m}^{r-1}h\otimes\omega)$ 与
$\beta_{c,m}(\ell_{\omega,\xi,m}^{r-1}h\otimes\omega)$ 都等于
\[
\omega^{\otimes r}\otimes\pi^{n-m}\kaTr_{m,n}(h(\xi_m))
(c\pi^{-n}\log(\chi_G)).
\]
\begin{statement}{引理 2.3.9}
存在非零整数 $c$，满足 $(r-1)!^{-1}\pi^{-nr}c\in\mathcal O_K$，使得对所有 $m\geq n$，
$\tau_{c,m}$ 消去 $J_{S_m}^{[r-1]}\otimes_{\kaOG}\kaOm$。
\end{statement}
\origpage{146}{141}
这由 2.3.3 及下面的引理得到。
\begin{statement}{引理 2.3.10}
设 $\omega$ 是 $\kaco$ 的一组基。存在非零整数 $c$，对所有 $m\geq n$，它都消去
\[
J_{S_m}^{[r-1]}/\bigl(J_{S_m}^{[r]}+\ell_{\omega,\xi,m}^{r-1}S_m\bigr).
\]
\end{statement}
为证明 2.3.10，先记下如下引理。
\begin{statement}{引理 2.3.11}
对 $i\geq0$，$J_{S_m}^{[i]}/J_{S_m}^{[i+1]}$ 是秩为 $1$ 的自由
$\mathcal O_{K_m}=S_m/J_{S_m}$ 模。若 $h\in J_{S_m}\subset S_m$ 满足
$h\bmod J_{S_m}^{[2]}$ 是 $J_{S_m}/J_{S_m}^{[2]}$ 的一组基，则对任意 $i\geq0$，
$h^{[i]}$ 是 $J_{S_m}^{[i]}/J_{S_m}^{[i+1]}$ 的一组基。
\end{statement}
2.3.11 的证明从略。
\begin{proof}[2.3.10 的证明]
由 2.3.11，只需证明存在非零整数，对所有 $m\geq n$，都消去
$J_{S_m}/(J_{S_m}^{[2]}+\ell_{\omega,\xi,m}S_m)$。
令 $I=\Ker(\kaOG\to\mathcal O_K)$ 为 $\kaOG$ 的增广理想。则
$I/I^2\simeq J_{S_0}/J_{S_0}^{[2]}$。因为 $\ell_\omega\bmod J_{S_0}^{[2]}$
是 $J_{S_0}/J_{S_0}^{[2]}$ 的一组基，可取 $I$ 的生成元 $h$ 使
$h\equiv\ell_\omega\pmod{J_{S_0}^{[2]}}$。对 $m\geq0$，令 $h_m=[\pi^m]^*h$。
则对 $m\geq1$，$g_m=h_mh_{m-1}^{-1}$ 是理想
$\Ker(\xi_m:\kaOG\to\mathcal O_{K_m})$ 的生成元。有同构
\[
g_m\kaOG/g_m^2\kaOG\xrightarrow{\sim}J_{S_m}/J_{S_m}^{[2]},
\]
将 $h_m$ 的类映为 $\ell_{\omega,\xi,m}$ 的类。因此对 $m\geq1$，
\begin{align*}
J_{S_m}/(J_{S_m}^{[2]}+\ell_{\omega,\xi,m}S_m)
&\simeq g_m\kaOG/(g_m^2\kaOG+h_m\kaOG)\\
&\simeq\kaOG/(g_m\kaOG+h_{m-1}\kaOG)\\
&\simeq[\pi^{m-1}]^*\bigl(\kaOG/(g_1\kaOG+h\kaOG)\bigr).
\end{align*}
由于 $g_1$ 与 $h$ 在 $\kaOG$ 中没有公共素因子，
$\kaOG/(g_1\kaOG+h\kaOG)$ 被某个非零整数消去；同一个整数对所有 $m$
都消去 $J_{S_m}/(J_{S_m}^{[2]}+\ell_{\omega,\xi,m}S_m)$。
\end{proof}

\origpage{147}{142}
\noindent\textbf{2.3.12.} 为从 2.3.9 推出命题 2.3.2，考虑 $a\in\mathcal O_F$ 的作用
\[
[a]^*:S_m\otimes_{\kaOG}\kaOm\longrightarrow S_m\otimes_{\kaOG}\kaOm.
\]
对任意满足 $(r-1)!^{-1}\pi^{-nr}c\in\mathcal O_K$ 的整数 $c$，以及任意
$a\in\Ker(\mathcal O_F^\times\to(\mathcal O_F/\pi^n)^\times)$，容易看出
\[
\alpha_{c,m}\circ[a]^*=a^r\alpha_{c,m},\qquad
\beta_{c,m}\circ[a]^*=a^r\beta_{c,m},
\]
从而 $\tau_{c,m}\circ[a]^*=a^r\tau_{c,m}$。由此，命题 2.3.2 可归结为 2.3.9 和下面的引理。
\begin{statement}{引理 2.3.13}
令 $C_m=(S_m/J_{S_m}^{[r-1]})\otimes_{\kaOG}\kaOm$，并令 $D_m\subset C_m$
为对所有 $a\in\Ker(\mathcal O_F^\times\to(\mathcal O_F/\pi^n)^\times)$ 的
$([a]^*-a^r)C_m$ 所生成的子群。则存在非零整数，对所有 $m\geq n$ 都消去 $C_m/D_m$。
\end{statement}
2.3.13 的证明可归结为下面的引理。
\begin{statement}{引理 2.3.14}
设 $1\leq i\leq r-1$。令 $E_m=(J_{S_m}^{[i-1]}/J_{S_m}^{[i]})\otimes_{\kaOG}\kaOm$，
并令 $F_m$ 为对所有 $a\in\Ker(\mathcal O_F^\times\to(\mathcal O_F/\pi^n)^\times)$ 的
$([a]^*-a^r)E_m$ 所生成的 $E_m$ 的子群。则存在非零整数，对所有 $m\geq n$ 都消去 $E_m/F_m$。
\end{statement}
\begin{proof}
取 2.3.10 证明中的 $g_m$，并取 $\kaco$ 的一组基 $\omega$，考虑
$\mathcal O_{K_m}$ 模的同构
\[
f:\mathcal O_{K_m}\xrightarrow{\sim}
(J_{S_m}^{[i-1]}/J_{S_m}^{[i]})\otimes_{\kaOG}\kaOm,
\quad1\longmapsto g_m^{[i-1]}\otimes\omega.
\]
则对于任意 $a\in\Ker(\mathcal O_F^\times\to(\mathcal O_F/\pi^n)^\times)$，
有 $fa^i\sigma_a=[a]^*f$，其中
$\sigma_a\in\Gal((\bigcup_jK_j)/K_n)$ 是在 $T$ 上作用为 $a$ 的唯一元素。
$F_m$（2.3.14）在 $f$ 下的原像，是对所有上述 $a$ 的
$(a^i\sigma_a-a^r)\mathcal O_{K_m}$ 所生成的 $\mathcal O_{K_m}$ 的子群。
\end{proof}
\origpage{148}{143}
因此归结为下面的引理。
\begin{statement}{引理 2.3.15}
设 $r$ 是非零整数，$A_m$ 为对所有
$a\in\Ker(\mathcal O_F^\times\to(\mathcal O_F/\pi^n)^\times)$ 的
$(\sigma_a-a^r)\mathcal O_{K_m}$ 所生成的 $\mathcal O_{K_m}$ 的子群。
则存在非零整数，对所有 $m\geq n$ 都消去 $\mathcal O_{K_m}/A_m$。
\end{statement}
\begin{proof}
因为 $A_m$ 是 $\mathcal O_{K_m}$ 的 $\mathcal O_{K_n}$ 子模，并对某个 $i\geq0$
包含 $p^i\mathcal O_{K_m}$，而且配对
\[
\mathcal O_{K_m}/p^i\times\mathcal O_{K_m}/p^i\longrightarrow\mathcal O_{K_n}/p^i,
\qquad(x,y)\longmapsto\pi^{n-m}\kaTr_{m,n}(xy)
\]
是有限生成自由 $\mathcal O_{K_n}/p^i$ 模之间的完美对偶，考虑对偶即可归结为下面的引理。
\end{proof}
\begin{statement}{引理 2.3.16}
固定 $r\in\mathbb Z$，$r\ne0$。存在非零整数，对所有 $m\geq n$ 及 $i\geq0$，都消去
\[
\left\{x\in\mathcal O_{K_m}/p^i\ ;\ 
\sigma_a(x)=a^rx\ \text{对一切 }a\in
\Ker(\mathcal O_F^\times\to(\mathcal O_F/\pi^n)^\times)\right\}.
\]
\end{statement}
\begin{proof}
最后这个群同构于
$H^0(K_n,(\mathcal O_{K_m}/p^i)\otimes_{\mathcal O_F}T^{\otimes(-r)})$，
后者嵌入
$H^0(K_n,(\mathcal O_{\bar K}/p^i)\otimes_{\mathcal O_F}T^{\otimes(-r)})$。
因此，只需证明
$H^q(K_n,\widehat{\mathcal O}_{\bar K}\otimes_{\mathcal O_F}T^{\otimes(-r)})$
（$q=0,1$）被 $p$ 的某个幂消去。因为
\[
(\widehat{\mathcal O}_{\bar K}\otimes_{\mathcal O_F}T^{\otimes(-r)})\otimes\mathbb Q
\simeq(\kaco^{\otimes r}\otimes_{\mathcal O_K}
\widehat{\mathcal O}_{\bar K}(-r))\otimes\mathbb Q,
\]
这就归结为 $H^q(K_n,\widehat{\mathcal O}_{\bar K}(-r))$ 被 $p$ 的某个幂消去这一事实
（2.2.4 (1)）。
\end{proof}

\origpage{149}{144}
\section*{第三章\quad 整体问题}
\addcontentsline{toc}{section}{第三章 整体问题}
\subsection*{§1. 显式互反律与 zeta 值}
\addcontentsline{toc}{subsection}{§1. 显式互反律与 zeta 值}
我们将第二章 §2 中的\term{显式互反律}{explicit reciprocity law}{loi de réciprocité explicite}
与部分 Riemann zeta 函数的值，以及类数为 $1$ 的虚二次域的某些 Hecke $L$ 级数的值联系起来
（定理 1.2.6）。§1.1 是预备性的回顾。

\subsection*{§1.1. $\mathbb G_m$ 及椭圆曲线上的基本函数与 zeta 值（回顾）}
在本 §1.1 中，设 $K$ 是域，$G$ 是 $K$ 上的乘法群 $\mathbb G_{m,K}$ 或 $K$ 上的一条
\term{椭圆曲线}{elliptic curve}{courbe elliptique}。以 $K(G)$ 表示 $G$ 上的有理函数域。

\noindent\textbf{1.1.1.} 对于 $a\in\operatorname{End}(G)\setminus\{0\}$，令
$a^*:K(G)\to K(G)$ 为沿 $a$ 的拉回，令 $\deg(a)$ 为这个域同态的次数，
$N_a:K(G)^\times\to K(G)^\times$ 为 $a^*$ 的范同态。
对 $f\in K(G)^\times$ 有 $N_a(a^*f)=f^{\deg(a)}$。
当 $a\in\mathbb Z$ 时，若 $G=\mathbb G_{m,K}$，则 $\deg(a)=|a|$；若 $G$ 是椭圆曲线，
则 $\deg(a)=a^2$。
\begin{statement}{定义 1.1.2}
若元素 $\theta\in K(G)^\times$ 满足以下两个条件，则称它为 $G$ 上的
\term{基本函数}{fundamental function}{fonction fondamentale}：
\begin{enumerate}[label=(\arabic*)]
\item 对任意 $a\in\operatorname{End}(G)\setminus\{0\}$，有 $N_a(\theta)=\theta$。
\item 除了在 $G$ 的原点处有一个一阶零点之外，$\theta$ 在 $G$ 上没有其他零点或极点。
\end{enumerate}
\end{statement}
$G$ 上的基本函数不存在。（当 $G$ 是椭圆曲线时，由主除子的次数为零，与条件 (2) 矛盾即可看出。）
\origpage{150}{145}
显然，作者应当解释为什么要写下这个空的定义 1.1.2。他认为，对于某些 zeta 值的研究，
从下面这句错误的话中找出“某种真意”很有帮助：“基本函数存在且唯一。当
$G=\mathbb G_{m,K}$ 或 $G$ 是具有复乘的椭圆曲线时，某些 zeta 值可以由基本函数的对数
反复求导后在 $G$ 的挠点处的取值得到。”

\noindent\textbf{1.1.3.} 对 $G=\mathbb G_{m,K}=\Spec(K[t,t^{-1}])$（$t$ 为
$\mathbb G_{m,K}$ 的标准坐标），函数 $(1-t)(1-t^{-1})$ 起着 $\theta^2$ 的作用。
事实上，它是 $K(G)^\times=K(t)^\times$ 中唯一对所有
$a\in\operatorname{End}(G)\setminus\{0\}=\mathbb Z\setminus\{0\}$ 都在 $N_a$ 下不变、
而且除了原点处的二阶零点外在 $G$ 上没有其他零点或极点的元素。
对于 $G=\mathbb G_{m,K}$，1.1.2 中的“某种真意”是指
\[
\left(t\frac d{dt}\right)^r\log\bigl((1-t)(1-t^{-1})\bigr)\quad(r\geq0)
\]
在 $\mathbb G_{m,K}$ 的挠点处的值，是部分 Riemann zeta 函数的值
（第一章 §1.3、第一章 3.3.6；也参见下文 1.1.7）。

\noindent\textbf{1.1.4.} 以 $e$ 表示 $G$ 的原点。对 $a,b\in\operatorname{End}(G)$，
若作为概形有 $\Ker(a)\cap\Ker(b)=\{e\}$（$\{e\}$ 赋予既约概形结构），且 $ab=ba$，
则记 $(a,b)=1$。（若 $G=\mathbb G_{m,K}$ 或 $\operatorname{char}(K)=0$，
$\operatorname{End}(G)$ 是交换的，因而 $ab=ba$ 自动成立。）下面的 1.1.5 表明：
若 $G$ 是椭圆曲线，$a\in\operatorname{End}(G)$ 满足 $(a,6)=1$，则存在函数
$\theta_a\in K(G)^\times$，起着 $\theta^{a^2}(a^*\theta)^{-1}$ 的作用。
\begin{statement}{命题 1.1.5}
设 $G$ 是椭圆曲线，$a\in\operatorname{End}(G)$，$(a,6)=1$。则存在唯一的有理函数
$\theta_a\in K(G)^\times$ 满足以下条件 (1)、(2)：
\origpage{151}{146}
\begin{enumerate}[label=(\arabic*)]
\item 对任意满足 $(a,b)=1$ 的 $b\in\operatorname{End}(G)$，有 $N_b(\theta_a)=\theta_a$。
\item $\theta_a$ 的除子为 $\deg(a)(e)-\Ker(a)$。
\end{enumerate}
对任意满足 $(b,6)=1$ 且 $ab=ba$ 的 $b\in\operatorname{End}(G)$，这个函数 $\theta_a$ 还满足
\begin{equation*}\tag{3}
(\theta_a)^{\deg(b)}(b^*\theta_a)^{-1}
=(\theta_b)^{\deg(a)}(a^*\theta_b)^{-1}.
\end{equation*}
\end{statement}
1.1.5 的证明将在下文 1.1.13 给出。

\noindent\textbf{1.1.6.} 现在设 $K$ 为\term{类数}{class number}{nombre de classes}为 $1$ 的
\term{虚二次域}{imaginary quadratic field}{corps quadratique imaginaire}，$E$ 为 $K$ 上的一条椭圆曲线，
使得 $\operatorname{End}(E)\to\operatorname{End}(\operatorname{Lie}(E))=K$ 诱导同构
$\operatorname{End}(E)\xrightarrow{\sim}\mathcal O_K$；我们把它视为恒等识别。则
\[
L(H^1(E),s)=\sum_{\mathfrak a}\overline{\nu(\mathfrak a)}N(\mathfrak a)^{-s}
=\sum_{\mathfrak a}\nu(\mathfrak a)^{-1}N(\mathfrak a)^{-(s-1)},
\]
其中 $\nu$ 是 $K$ 的一个\term{Hecke特征}{Hecke character}{caractère de Hecke}，满足：
对于 $\mathcal O_K$ 的任意与 $\nu$ 的\term{导子}{conductor}{conducteur} $\mathfrak f$ 互素的非零理想
$\mathfrak a$，$\nu(\mathfrak a)$ 是 $\mathfrak a$ 的一个生成元；$\overline{\phantom x}$ 表示复共轭，
求和中的 $\mathfrak a$ 遍历 $\mathcal O_K$ 的所有与 $\mathfrak f$ 互素的理想。

设 $\mathfrak g$ 是 $\mathcal O_K$ 的非零理想，$\mathfrak g\subset\mathfrak f$；令 $L$ 为 $K$ 的如下阿贝尔扩张：
它对应于映射
\[
\Gal(\bar K/K)\longrightarrow\operatorname{Aut}_{\mathcal O_K}({}_{\mathfrak g}E)
\simeq(\mathcal O_K/\mathfrak g)^\times
\]
的核。这里 ${}_{\mathfrak g}E$ 表示 $E$ 的 $\mathfrak g$ 挠点子群概形。
对 $r\in\mathbb Z$、$r\geq0$ 及 $\sigma\in\Gal(L/K)$，我们考虑部分 $L$ 函数
\[
L(\nu^{-r},\sigma\text{-part},s)
=\sum_{(\mathfrak a,L/K)=\sigma}\nu(\mathfrak a)^{-r}N(\mathfrak a)^{-s},
\]
求和遍历 $\mathcal O_K$ 的所有与 $\mathfrak g$ 互素、且
\term{Artin符号}{Artin symbol}{symbole d'Artin}
$(\mathfrak a,L/K)\in\Gal(L/K)$ 等于 $\sigma$ 的非零理想 $\mathfrak a$。
这个 $L$ 函数在 $\Re(s)>1-r/2$ 时绝对收敛。
\origpage{152}{147}
当 $r\geq1$（分别地，$r=0$）时，它可延拓为整个复平面上的全纯（分别地，亚纯）函数。
固定嵌入 $K\hookrightarrow\mathbb C$。令
\[
\exp:\operatorname{Lie}(E)\otimes_K\mathbb C\longrightarrow E(\mathbb C)
\]
为指数映射。以下定理是经典结果（参见 [Da]、[We1]、[dS2]）；其中 (2) 是 Kronecker
的“\term{第二极限公式}{second limit formula}{seconde formule limite}”。
\begin{statement}{定理}
设 $\gamma$ 为可逆 $\mathcal O_K$ 模 $\mathfrak g^{-1}H_1(E(\mathbb C),\mathbb Z)$ 的一个生成元，
$\alpha=\exp(\gamma)\in{}_{\mathfrak g}E(\mathbb C)={}_{\mathfrak g}E(L)$。
设 $\mathfrak a$ 是 $\mathcal O_K$ 的与 $6\mathfrak g$ 互素的非零理想，
$\sigma\in\Gal(L/K)$，$\tau\overset{\mathrm{def}}=(\mathfrak a,L/K)\in\Gal(L/K)$。
\begin{enumerate}[label=(\arabic*)]
\item 设 $r\in\mathbb Z$、$r\geq1$，$\omega\in\operatorname{coLie}(E)\setminus\{0\}$。
则 $(d/\omega)^r\log(\theta_{\nu(\mathfrak a)})$ 在 $\sigma\alpha$ 处的值等于
\[
(-1)^{r-1}(r-1)!\left(\int_\gamma\omega\right)^{-r}
\]
乘以
\[
N(\mathfrak a)L(\nu^{-r},\sigma\text{-part},0)
-\nu(\mathfrak a)^rL(\nu^{-r},\sigma\tau\text{-part},0).
\]
\item 有
\begin{align*}
\log\left|\theta_{\nu(\mathfrak a)}(\sigma\alpha)\right|
=-\lim_{s\to0}s^{-1}\bigl(&N(\mathfrak a)L(1,\sigma\text{-part},s)\\
&-L(1,\sigma\tau\text{-part},s)\bigr).
\end{align*}
\end{enumerate}
\end{statement}
\noindent\textbf{1.1.7.} 上述定理类似于第一章 §1.3 和 3.3.6 中引入的部分 Riemann zeta 函数的下列性质。

设 $\alpha=\exp(2\pi iaN^{-1})$，其中 $N,a\in\mathbb Z$、$N\geq1$，并假设 $\alpha\ne1$。则：
\begin{enumerate}[label=(\arabic*)]
\item 对 $r\in\mathbb Z$、$r\geq1$，函数
$\tfrac12(t\frac d{dt})^r\log((1-t)(1-t^{-1}))$ 在 $\alpha$ 处的值等于
\[
(-1)^{r-1}(r-1)!\left(\frac{2\pi i}{N}\right)^{-r}
\bigl(\zeta_{a(N)}+(-1)^r\zeta_{-a(N)}\bigr)(r).
\]
\item 有
\[
\frac12\log((1-\alpha)(1-\alpha^{-1}))
=-\lim_{s\to0}s^{-1}\bigl(\zeta_{a(N)}(s)+\zeta_{-a(N)}(s)\bigr).
\]
\end{enumerate}

\origpage{153}{148}
无论在 1.1.6 还是 1.1.7 中，每挥动一次魔杖 $d/\omega$，就出现一次 zeta 值。
对于任意\term{动机}{motif}{motif}，是否也有这样一根魔杖？

\noindent\textbf{1.1.8.} 简要回顾 1.1.6 中定理 (1) 在 $r\geq3$ 时的证明。
（借助 $r\geq3$ 的假设，可以避开有关收敛性与解析延拓的精细论证。1.1.6 中定理的完整证明见 [We1]。）
我们使用下面的 1.1.9。

\noindent\textbf{1.1.9.} 设 $G$ 是 $\mathbb C$ 上的椭圆曲线，
$a\in\operatorname{End}(G)\setminus\{0\}$，$(a,6)=1$，
$r\in\mathbb Z$、$r\geq3$，$\omega\in\kaco\setminus\{0\}$。则复合亚纯函数
\[
\operatorname{Lie}(G)\xrightarrow{\exp}G(\mathbb C)\xrightarrow{(*)}\mathbb C,
\qquad(*)=\left(\frac d\omega\right)^r\log(\theta_a)
\]
等于
\[
z\longmapsto(-1)^{r-1}(r-1)!
\bigl(\deg(a)E_{r,\omega}(z)-\operatorname{Lie}(a)^rE_{r,\omega}(az)\bigr),
\]
其中 $\operatorname{Lie}(a)\in\mathbb C$ 是 $a$ 在 $\operatorname{Lie}(G)$ 上的作用，且
\[
E_{r,\omega}(z)=\sum_\gamma(z+\gamma,\omega)^{-r};
\]
这里 $\gamma$ 遍历
\[
H_1(G(\mathbb C),\mathbb Z)=\Ker(\exp:\operatorname{Lie}(G)\to G(\mathbb C)),
\]
而 $(\ ,\ )$ 是典范配对 $\operatorname{Lie}(G)\times\kaco\to\mathbb C$。

\noindent\textbf{1.1.10.} 1.1.9 类似于下述事实：对 $r\in\mathbb Z$、$r\geq1$，复合亚纯函数
\[
\mathbb C\xrightarrow{\exp}\mathbb C^\times\xrightarrow{(*)}\mathbb C,
\qquad(*)=\frac12\left(t\frac d{dt}\right)^r\log((1-t)(1-t^{-1}))
\]
等于
\[
\begin{cases}
z\longmapsto(-1)^{r-1}(r-1)!\displaystyle\sum_{n\in\mathbb Z}(z-2\pi in)^{-r},&r\geq2,\\[1ex]
z\longmapsto\frac12(-1)^{r-1}(r-1)!\displaystyle\sum_{n\in\mathbb Z}
\bigl((z-2\pi in)^{-1}+(z+2\pi in)^{-1}\bigr),&r=1.
\end{cases}
\]
\noindent\textbf{1.1.11.} 由 1.1.9 推出 1.1.6 中定理 (1) 在 $r\geq3$ 时的结论。
只需证明
\origpage{154}{149}
\[
E_{r,\omega}(\sigma\alpha)=\left(\int_\gamma\omega\right)^{-r}
L(\nu^{-r},\sigma\text{-part},0).
\]
这里使用了以下事实：对于 $a\in\mathcal O_K=\operatorname{End}(E)$ 及
$\mathfrak a=(a)$，有 $\deg(a)=a\bar a=N(\mathfrak a)$。
取 $b\in\mathcal O_K$，使 $\sigma\alpha=b\alpha$。则
\[
E_{r,\omega}(\sigma\alpha)=\sum_\delta(b\gamma+\delta,\omega)^{-r}
=\left(\sum_{x\in\mathfrak g}(b+x)^{-r}\right)
\left(\int_\gamma\omega\right)^{-r},
\]
其中 $\delta$ 遍历 $H_1(G(\mathbb C),\mathbb Z)$ 的所有元素。但是
\[
\sum_{x\in\mathfrak g}(b+x)^{-r}
=\sum_{(\mathfrak a,L/K)=\sigma}\nu(\mathfrak a)^{-r}
=L(\nu^{-r},\sigma\text{-part},0).
\]
第一个等式来自如下事实：映射
\[
\mathfrak a\longmapsto\nu(\mathfrak a),\qquad(a)\longleftarrow a
\]
给出互逆的双射，一边是 $\mathcal O_K$ 的所有与 $\mathfrak g$ 互素、且
$(\mathfrak a,L/K)=\sigma$ 的理想组成的集合，另一边是 $\mathcal O_K$ 中所有模
$\mathfrak g$ 同余于 $b$ 的元素组成的集合。

\noindent\textbf{1.1.12.} 类似地，由 1.1.10 可推出 1.1.7，并得到第一章 1.3.2 (2) 的另一个证明。

\noindent\textbf{1.1.13.} 现证明 1.1.5。由 Abel 定理，$\deg(a)(e)-\Ker(a)$ 是主除子。
任取 $g\in K(G)^\times$，使 $\operatorname{div}(g)=\deg(a)(e)-\Ker(a)$。
则对任意满足 $(a,b)=1$ 的 $b\in\operatorname{End}(G)$，有
$\operatorname{div}(N_b(g))=\operatorname{div}(g)$，因而存在 $\lambda_b\in K^\times$，
使 $N_b(g)=\lambda_bg$。令 $\theta_a=\lambda_3\lambda_2^{-3}g$。
我们证明 $\theta_a$ 满足条件 (1)、(2)。(2) 显然，因此证明 (1)。
对任意满足 $(a,b)=1$、$(a,c)=1$ 及 $bc=cb$ 的 $b,c\in\operatorname{End}(G)$，
由 $N_bN_c=N_cN_b$ 得
\[
\lambda_b\lambda_c^{\deg(b)}=\lambda_c\lambda_b^{\deg(c)}.
\]
考虑此式中 $c=2,3$ 的情形，得到：对任意满足 $(a,b)=1$ 的
$b\in\operatorname{End}(G)$，
$\lambda_b=(\lambda_2^3\lambda_3^{-1})^{\deg(b)-1}$。由此即知 (1) 成立。

$\theta_a$ 的唯一性如下证明。满足条件 (1)、(2) 的两个函数之比 $\alpha$，由 (1) 具有性质
\origpage{155}{150}
$\alpha=N_b(\alpha)=\alpha^{\deg(b)}$，其中 $b$ 任意满足 (1) 的条件。
取 $b=2,3$ 即得 $\alpha=1$。

最后证明 (3)。(3) 两边有相同的除子
\[
\deg(a)\deg(b)(e)-\deg(a)\Ker(b)-\deg(b)\Ker(a)+\Ker(ab).
\]
因此两边之比为 $K^\times$ 中的非零元素。对 (3) 两边分别施以 $N_2$ 和 $N_3$，便知此比为 $1$。

\subsection*{§1.2. 与显式互反律的关系}
\noindent\textbf{1.2.1.} 在本 §1.2 中，设 $p$ 为素数，$K$ 为以下两者之一：
\begin{enumerate}[label=(\arabic*)]
\item $\mathbb Q$；
\item 类数为 $1$ 的虚二次域。
\end{enumerate}
设 $S$ 是 $\mathcal O_K$ 的非零素理想的有限集合，包含 $p$ 的所有素因子，
并令 $X=\Spec(\mathcal O_K[1/S])$。
在第一章 §3.3 中，对于上述情形 (1)，我们尝试为第一章 3.1.1 中的 $\Lambda$ 及
$X$ 上的光滑可逆 $\Lambda$ 层 $\mathscr F$ 定义
\term{zeta元素}{zeta element}{élément zêta} $z_\Lambda(X,\mathscr F)$。
在情形 (2) 中也有类似的理论（1.2.3）。正如\term{分圆单位}{cyclotomic unit}{unité cyclotomique}
在情形 (1) 中起着关键作用，“\term{椭圆单位}{elliptic unit}{unité elliptique}”在情形 (2) 中起着关键作用。
我们将说明，无论在 (1) 还是 (2) 中，我们的 zeta 元素通常都借助第二章 §2 的显式互反律
与 zeta 值相联系，这也是第一章猜想 3.2.2 的条件 (v) 所预期的。

在情形 (2) 中，本 §1.2 固定 $K$ 上的一条椭圆曲线 $E$，满足
$\operatorname{End}(E)\simeq\mathcal O_K$。假设 $E$ 在 $S$ 之外都有好约化，
并且在 $p$ 的素因子处也有好约化。把 $\mathcal O_K$ 与 $\operatorname{End}(E)$ 识别，
使复合映射
\[
\mathcal O_K\longrightarrow\operatorname{End}(E)\longrightarrow
\operatorname{End}_K(\operatorname{Lie}(E))=K
\]
是自然包含。在情形 (2) 中，还固定嵌入 $K\hookrightarrow\mathbb C$。
\origpage{156}{151}
在情形 (1)、(2) 中都固定 $K$ 的一个代数闭包 $\bar K$，并令 $K^{\mathrm{ab},S}$ 为
$\bar K$ 中所有在 $S$ 之外不分歧的有限阿贝尔扩张的并。
本 §1.2 的写作受到了 Soulé [So4] 对“$p$ 进 Galois 上同调中的椭圆元素”的定义的影响。

\noindent\textbf{1.2.2.} 设处于 1.2.1 (2) 的情形。考虑第一章 3.3.5 的类似构造。令
\[
\Lambda_0=\mathbb Z_p[[\Gal(K^{\mathrm{ab},S}/K)]],\qquad
\mathscr F_0=\varprojlim_L(f_L)_*f_L^*(\mathbb Z_p(1)_X),
\]
其中 $L$ 遍历 $K^{\mathrm{ab},S}/K$ 的有限子扩张，$f_L$ 表示典范态射
$\Spec(\mathcal O_L[1/S])\to\Spec(\mathcal O_K[1/S])$。
取 $\mathcal O_K$ 的非零元素 $a$，与 $6$ 及 $S$ 的任意元素互素。
在本 1.2.2 中，定义元素
\[
z_{a,\Lambda_0}(X,\mathscr F_0)\in H^1(X,\mathscr F_0)
\otimes_{\Lambda_0}H^0(\mathbb C,\mathscr F_0(-1))^{-1}.
\]
$\theta_a$（§1.1）在 $E$ 的挠点处的值，即所谓椭圆单位，在此定义中起着关键作用。
这类似于 $(1-t)(1-t^{-1})$ 在 $\mathbb G_{m,K}=\Spec(K[t,t^{-1}])$ 的挠点处的值，
即所谓分圆单位，在第一章 3.3.5 中定义 $z_{\Lambda_+}(X,\mathscr F_+)$ 时起着关键作用。

设 $\mathfrak f$ 是与 $E$ 相伴的 Hecke 特征的导子（§1.1）。
对于满足 $\mathfrak g\subset\mathfrak f$，并且素因子集合恰为 $S$ 的
$\mathcal O_K$ 的非零理想 $\mathfrak g$，令 $K(\mathfrak g)\subset K^{\mathrm{ab},S}$ 为
$K$ 的有限阿贝尔扩张，对应于
$\Gal(K^{\mathrm{ab},S}/K)\to\operatorname{Aut}({}_{\mathfrak g}E)$ 的核，
其中 ${}_{\mathfrak g}E$ 是 $E$ 的 $\mathfrak g$ 挠点子群概形。
则 $K^{\mathrm{ab},S}=\bigcup_{\mathfrak g}K(\mathfrak g)$。令
\begin{equation*}\tag{$(1)_{a,\mathfrak g}$}
H^0(K(\mathfrak g)\otimes_K\mathbb C,\mathbb Z_p)
\simeq\bigoplus_\iota H^0(\mathbb C,\mathbb Z_p)
\longrightarrow H^1(\mathcal O_{K(\mathfrak g)}[1/S],\mathbb Z_p(1))
\end{equation*}
\origpage{157}{152}
为同态 $(a_\iota)_\iota\mapsto\sum_\iota a_\iota\{\iota^{-1}(\theta_a(x_{\mathfrak g}))\}$，
其中 $\iota$ 遍历所有限制到 $K$ 后与给定嵌入 $K\hookrightarrow\mathbb C$ 一致的
$K(\mathfrak g)\hookrightarrow\mathbb C$；
$x_{\mathfrak g}=\exp(\gamma_{\mathfrak g})$，$\gamma_{\mathfrak g}$ 是
$\mathfrak g^{-1}H_1(E(\mathbb C),\mathbb Z)$ 的一组 $\mathcal O_K$ 基。
此时 $\theta_a(x_{\mathfrak g})$ 与 $\gamma_{\mathfrak g}$ 的选择无关，因为由 $\theta_a$ 的性质
1.1.5 (1)，对任意 $u\in\mathcal O_K^\times$ 都有 $\theta_a\circ[u]=\theta_a$。
已知 $\theta_a(x_{\mathfrak g})\in(\mathcal O_{K(\mathfrak g)}[1/S])^\times$
（[dS2] 第二章 §2）。映射 $(1)_{a,\mathfrak g}$ 随 $\mathfrak g$ 变化构成一个逆系统。
其逆极限
\[
H^0(\mathbb C,\mathscr F_0(-1))\longrightarrow H^1(X,\mathscr F_0)
\]
对应于 $H^1(X,\mathscr F_0)\otimes_{\Lambda_0}H^0(\mathbb C,\mathscr F_0(-1))^{-1}$
中的一个元素，将它记为 $z_{a,\Lambda_0}(X,\mathscr F_0)$。

\noindent\textbf{1.2.3.} 设处于 1.2.1 (2) 的情形。考虑第一章 3.3.7 的类似构造。
对第一章 3.1.1 中的任意 $\Lambda$ 及 $X$ 上任意光滑可逆 $\Lambda$ 层 $\mathscr F$，
尝试定义 $\Delta_\Lambda(X,\mathscr F)$ 的一个典范 $\Lambda$ 基 $z_\Lambda(X,\mathscr F)$。
定义\footnote{译注：原书152--154页的若干张量式在 $H^0$ 因子后未印逆模指数 $-1$，而前后定义有此指数，疑为漏印；此处及后文均照录原式。}
\[
z_{\Lambda_0}(X,\mathscr F_0)\in H^1(X,\mathscr F_0)\otimes_{\Lambda_0}
H^0(\mathbb C,\mathscr F_0(-1))\otimes_{\Lambda_0}Q(\Lambda_0)
\]
为
\[
z_{\Lambda_0}(X,\mathscr F_0)
=(N(\mathfrak a)-\sigma_{\mathfrak a})^{-1}
z_{\nu(\mathfrak a),\Lambda_0}(X,\mathscr F_0),
\]
其中 $\mathfrak a$ 为 $\mathcal O_K$ 的理想，满足 $(\mathfrak a,6)=1$、
$\mathfrak a\ne\mathcal O_K$，且与 $S$ 的每个元素互素。
则 $z_{\Lambda_0}(X,\mathscr F_0)$ 与 $\mathfrak a$ 的选择无关。
我们希望将它视为 $\Delta_{\Lambda_0}(X,\mathscr F_0)$ 的一组 $\Lambda_0$ 基。
一旦做到这一点，便可用第一章 3.3.2 中相同的方法，对任意上述 $(\Lambda,\mathscr F)$
得到 $z_\Lambda(X,\mathscr F)$。问题可归结为下面这个类似于第一章 3.3.8 的猜想。
设 $\mathfrak g$ 为 $\mathcal O_K$ 的理想，使 $\mathfrak gp$ 的素因子集合恰为 $S$，
令
\[
\begin{gathered}
\Lambda=\mathbb Z_p[[G_{\mathfrak gp^\infty}]],\qquad
G_{\mathfrak gp^\infty}=\varprojlim_n\Gal(K(\mathfrak gp^n)/K),\\
\mathscr F_{\mathfrak gp^\infty}=\mathscr F_0\otimes_{\Lambda_0}\Lambda_{\mathfrak gp^\infty}.
\end{gathered}
\]
出于与第一章 3.3.7 相同的理由，
\origpage{158}{153}
$H^i(X,\mathscr F_{\mathfrak gp^\infty})\otimes_{\Lambda_{\mathfrak gp^\infty}}
Q(\Lambda_{\mathfrak gp^\infty})$ 在 $i\ne1$ 时为零，而在 $i=1$ 时是可逆的
$Q(\Lambda_{\mathfrak gp^\infty})$ 模。
\begin{statement}{猜想}
$z_{\Lambda_0}(X,\mathscr F_0)$ 在
\begin{align*}
&\Delta_{\Lambda_{\mathfrak gp^\infty}}(X,\mathscr F_{\mathfrak gp^\infty})
\otimes_{\Lambda_{\mathfrak gp^\infty}}Q(\Lambda_{\mathfrak gp^\infty})\\
&\quad=H^1(X,\mathscr F_{\mathfrak gp^\infty})\otimes_{\Lambda_{\mathfrak gp^\infty}}
H^0(\mathbb C,\mathscr F_{\mathfrak gp^\infty}(-1))
\otimes_{\Lambda_{\mathfrak gp^\infty}}Q(\Lambda_{\mathfrak gp^\infty})
\end{align*}
中的像，是 $\Delta_{\Lambda_{\mathfrak gp^\infty}}(X,\mathscr F_{\mathfrak gp^\infty})$
的一组 $\Lambda_{\mathfrak gp^\infty}$ 基。
\end{statement}
对于多种 $S$，这个猜想可由 Rubin [Ru] 证明的“虚二次域 Iwasawa 理论的主猜想”的解得到。
上述猜想与“主猜想”之间的微妙区别，见第一章 3.3.9。

\noindent\textbf{1.2.4.} 下面的 1.2.4、1.2.5 为定理 1.2.6 作准备。
设 $A$ 是 $K$ 的有限扩张的有限直积，$\Lambda=A\otimes\mathbb Q_p$。
设 $\mathscr F$ 是 $X$ 上的可逆光滑 $\Lambda$ 层，使得存在整数 $r\geq1$ 及同态
$\lambda:\Gal(K^{\mathrm{ab},S}/K)\to A^\times$，它对 $A^\times$ 的离散拓扑连续，
并满足如下条件。在情形 1.2.1 (1)（分别地，(2)）中，
\[
\chi_{\mathscr F}(\sigma)=\chi_{\mathrm{cyclo}}(\sigma)^{1-r}\lambda(\sigma)
\quad\text{对一切 }\sigma\in\Gal(\mathbb Q^{\mathrm{ab},S}/\mathbb Q),
\]
并且 $\chi_{\mathscr F}(\tau)=-1$，其中
$\tau\in\Gal(\mathbb Q^{\mathrm{ab},S}/\mathbb Q)$ 为复共轭；
分别地，
\[
\chi_{\mathscr F}(\sigma)=\chi_{\mathrm{cyclo}}(\sigma)\chi_E(\sigma)^{-r}\lambda(\sigma)
\quad\text{对一切 }\sigma\in\Gal(K^{\mathrm{ab},S}/K),
\]
其中 $\chi_E:\Gal(K^{\mathrm{ab},S}/K)\to(\mathcal O_K\otimes\mathbb Z_p)^\times
\subset\Lambda^\times$ 是在 $p$ 进 Tate 模 $T_pE$ 上的作用。
我们定义元素
\[
z_\Lambda(X,\mathscr F)\in H^1(X,\mathscr F)
\otimes_\Lambda H^0(K\otimes\mathbb R,\mathscr F(-1))^{-1},
\]
\origpage{159}{154}
其定义使用第一章 3.3.5 中的元素 $z_{\Lambda_+}(X,\mathscr F_+)$；那里假设了 $p\ne2$，
但本节允许 $p=2$。分别地，使用 1.2.2 中的 $z_{a,\Lambda_0}(X,\mathscr F_0)$，
其中 $a\in\mathcal O_K$ 与 $6$ 及 $S$ 的每个元素互素。
分别以 $\Lambda'$、$\mathscr F'$ 表示 $\Lambda_+$、$\mathscr F_+$（分别地，$\Lambda_0$、$\mathscr F_0$），
则有一个非典范的 $\Lambda$ 同构
\begin{equation*}\tag{1}
\mathscr F\simeq\mathscr F'\otimes_{\Lambda'}\Lambda,
\qquad\Lambda'\longrightarrow\Lambda:\sigma\longmapsto
\chi_{\mathrm{cyclo}}(\sigma)\chi_{\mathscr F}(\sigma)^{-1}.
\end{equation*}
由 (1) 诱导同态
\begin{equation*}\tag{2}
\begin{gathered}
H^1(X,\mathscr F')\otimes_{\Lambda'}H^0(K\otimes\mathbb R,\mathscr F'(-1))^{-1}\\
\longrightarrow H^1(X,\mathscr F)\otimes_\Lambda H^0(K\otimes\mathbb R,\mathscr F(-1))^{-1},
\end{gathered}
\end{equation*}
它与 $\Lambda$ 同构 (1) 的选择无关。

情形 1.2.1 (1)，$p\ne2$：定义 $z_\Lambda(X,\mathscr F)$ 为
$z_{\Lambda_+}(X,\mathscr F_+)$ 在映射 (2) 下的像。

情形 1.2.1 (1)，$p=2$：用与 $p\ne2$ 时相同的方法定义
\[
z_{\Lambda_+}(X,\mathscr F_+)\in H^1(X,\mathscr F_+)
\otimes H^0(K\otimes\mathbb R,\mathscr F_+(-1))\otimes\mathbb Q.
\]
虽然 $z_+(X,\mathscr F_+)$ 的定义中出现了 $2^{-1}$（第一章 3.3.5），但张量乘上 $\mathbb Q$ 后
这不再构成问题。定义 $z_\Lambda(X,\mathscr F)$ 为 $z_{\Lambda_+}(X,\mathscr F_+)$ 在映射 (2) 下的像。

情形 1.2.1 (2)：定义 $z_{a,\Lambda}(X,\mathscr F)$ 为
$z_{a,\Lambda_0}(X,\mathscr F_0)$ 在映射 (2) 下的像。
取 $\mathcal O_K$ 的非零理想 $\mathfrak a$，使 $(\mathfrak a,6\mathfrak g)=1$\footnote{译注：$6\mathfrak g$ 原书如此；本段未重新指定 $\mathfrak g$，此处照录。}，且
$N(\mathfrak a)\nu(\mathfrak a)^{-r}$ 不是 $1$ 的根。
令 $\sigma_{\mathfrak a}=(\mathfrak a,K^{\mathrm{ab},S}/K)$ 为 Artin 符号，并定义
\[
z_\Lambda(X,\mathscr F)=
\bigl(N(\mathfrak a)-N(\mathfrak a)\chi_{\mathscr F}(\sigma_{\mathfrak a})^{-1}\bigr)^{-1}
z_{\nu(\mathfrak a),\Lambda}(X,\mathscr F).
\]
括号内的元素在 $\Lambda$ 中可逆，因为
$\chi_{\mathscr F}(\sigma_{\mathfrak a})=N(\mathfrak a)\nu(\mathfrak a)^{-r}\alpha$，
其中 $\alpha$ 是 $\Lambda$ 中 $1$ 的某个根。于是 $z_\Lambda(X,\mathscr F)$ 与
\origpage{160}{155}
$\mathfrak a$ 的选择无关。

\noindent\textbf{1.2.5.} 设 $A$、$\mathscr F$ 如 1.2.4。如下定义可逆 $A$ 模 $\Sigma$ 和 $\Omega$。
取 $K^{\mathrm{ab},S}/K$ 的有限子扩张 $L/K$，使 $\lambda$（1.2.4）经由 $\Gal(L/K)$ 分解。
在情形 1.2.1 (1)（分别地，(2)）中，令
\[
\Sigma'=\bigoplus_\iota\mathbb Q(2\pi i)^{-r},\qquad\Omega'=L
\]
（分别地，$\Sigma'=\bigoplus_\iota H_1(E(\mathbb C),\mathbb Q)^{\otimes(-r)}$，
$\Omega'=\operatorname{coLie}(E)^{\otimes r}\otimes_KL$），
其中 $r$ 如 1.2.4，$\iota$ 遍历所有嵌入 $L\hookrightarrow\mathbb C$
（分别地，遍历所有限制到 $K$ 后与给定嵌入 $K\hookrightarrow\mathbb C$ 一致的嵌入，
且 $\otimes(-r)$ 作为可逆 $K$ 模的张量幂）。令
\[
\Sigma=\Sigma'\otimes_{K[\Gal(L/K)]}A,\qquad
\Omega=\Omega'\otimes_{K[\Gal(L/K)]}A.
\]
有 $\Lambda$ 模的典范同构\translatornote{15}
\[
\Sigma\otimes\mathbb Q_p\simeq H^0(K\otimes\mathbb R,\mathscr F(-1)),
\qquad\Omega\otimes\mathbb Q_p\simeq D_{\mathrm{dR}}(\mathscr F).
\]
定义 $A\otimes\mathbb R$ 模的同构
\[
a:\Omega\otimes\mathbb R\xrightarrow{\sim}\Sigma\otimes\mathbb R,
\]
或等价地，$A\otimes\mathbb R$ 模的同构
$(\Omega\otimes_A\Sigma^{-1})\otimes\mathbb R\simeq A\otimes\mathbb R$。
令 $a$ 为由如下定义的 $a':\Omega'\otimes\mathbb R\xrightarrow{\sim}\Sigma'\otimes\mathbb R$\translatornote{16}
诱导的映射。在情形 1.2.1 (1) 中，令 $a'$ 为复合映射
\[
\begin{gathered}
\Omega'\otimes\mathbb R=L\otimes\mathbb R\longrightarrow\bigoplus_\iota\mathbb C
\longrightarrow\bigoplus_\iota\mathbb C/(\mathbb R(2\pi i)^{1-r})\\
\xleftarrow{\sim}\bigoplus_\iota\mathbb R(2\pi i)^{-r}=\Sigma'\otimes\mathbb R,
\end{gathered}
\]
其中第一个箭头由嵌入 $\iota:L\hookrightarrow\mathbb C$ 诱导。
在情形 1.2.1 (2) 中，令 $a'$ 为复合映射
\[
\Omega'\otimes\mathbb R=\operatorname{coLie}(E)^{\otimes r}\otimes_KL\otimes_{\mathbb Q}\mathbb R
\longrightarrow\bigoplus_\iota\operatorname{coLie}(E)^{\otimes r}\otimes_K\mathbb C
\]
\origpage{161}{156}
\[
\xrightarrow{\sim}\bigoplus_\iota\Hom_K(H_1(E(\mathbb C),\mathbb Q)^{\otimes r},\mathbb C)
=\bigoplus_\iota H_1(E(\mathbb C),\mathbb Q)^{\otimes(-r)}\otimes_K\mathbb C
\simeq\Sigma'\otimes\mathbb R,
\]
其中前一行的第二个箭头由嵌入 $\iota:L\hookrightarrow\mathbb C$ 诱导，
后一行的第一个箭头由下式给出：
\[
\omega^{\otimes r}\longmapsto
\left(\gamma^{\otimes r}\longmapsto\left(\int_\gamma\omega\right)^r\right)
\quad\bigl(\omega\in\operatorname{coLie}(E)\setminus\{0\},\ 
\gamma\in H_1(E(\mathbb C),\mathbb Q)\setminus\{0\}\bigr).
\]
下面的定理 1.2.6 使用记号 $z_\Lambda(X,\mathscr F)$、$\Sigma$、$\Omega$。
这些记号在前文有时是在假设某些猜想成立后使用的。然而，1.2.6 中的
$z_\Lambda(X,\mathscr F)$（分别地，$\Sigma$、$\Omega$）已在 1.2.4（分别地，1.2.5）
中不依赖任何猜想地定义；定理 1.2.6 也不假设任何猜想成立。
\begin{statement}{定理 1.2.6}
设 $A$、$\mathscr F$ 及 $r$ 如 1.2.4。则 1.2.4 中定义的元素
$z_\Lambda(X,\mathscr F)$ 经映射
\[
\exp^*\otimes\id:H^1(X,\mathscr F)\otimes_\Lambda\Sigma^{-1}
\longrightarrow(\Omega\otimes_A\Sigma^{-1})\otimes\mathbb Q_p
\]
映到 $\Omega\otimes_A\Sigma^{-1}$ 中的一个元素，而该元素在 $A\otimes\mathbb R$ 中的像
（1.2.5）等于
\[
(-1)^{r-1}L_{A,S}(\mathscr F^*(1),0).
\]
$\Sigma$、$\Omega$ 的记号参见 1.2.5。
\end{statement}
1.2.6 的 $r=1$ 情形是已知结果，在情形 1.2.1 (2) 下，它与 Coates--Wiles [CW] 的工作相关，
见 [dS2] 第四章 §2 的说明。[BK] (7.13) 已得到 1.2.6 的一部分。
我们在此给出情形 1.2.1 (2) 下的证明，省略类似而更容易的情形 1.2.1 (1) 的证明。

设 $\mathfrak p$ 为 $\mathcal O_K$ 中位于 $p$ 上方的素理想，
$K_{\mathfrak p}$ 为 $K$ 在 $\mathfrak p$ 处的局部域，$\mathcal O_{\mathfrak p}$
为其赋值环。由于已经假设（1.2.1）$E$ 在 $\mathfrak p$ 处有好约化，存在
$\mathcal O_{\mathfrak p}$ 上的椭圆曲线 $\mathcal E$，满足
$\mathcal E\otimes_{\mathcal O_{\mathfrak p}}K_{\mathfrak p}=E\otimes_KK_{\mathfrak p}$。
于是 $\mathcal E$ 的形式完备化可识别为 $\mathcal O_{\mathfrak p}$ 上的
\term{Lubin--Tate群}{Lubin--Tate group}{groupe de Lubin--Tate}，
\origpage{162}{157}
它对应于二元组 $(K_{\mathfrak p},\pi)$，其中 $\pi=\nu(\mathfrak p)$。
设 $\mathfrak h$ 为 $\mathcal O_K$ 的理想，满足 $\mathfrak h\subset\mathfrak f$，
且 $\mathfrak h$ 的素因子集合为 $S\setminus\{\mathfrak p\}$。
令 $K'_{\mathfrak p}$ 为合成域 $K(\mathfrak h)K_{\mathfrak p}$，并把
$(K_{\mathfrak p},K'_{\mathfrak p})$ 看作第二章 §2 中的 $(F,K)$。则对任意 $n\geq0$，有
\[
(K'_{\mathfrak p})_n=K(\mathfrak h\mathfrak p^n)K_{\mathfrak p}.
\]
现在取 $H_1(E(\mathbb C),\mathbb Z)$ 的一组 $\mathcal O_K$ 基 $\delta$。
另取嵌入 $K^{\mathrm{ab},S}\hookrightarrow\mathbb C$，以 $\iota_n$ 表示它限制到
$K(\mathfrak h\mathfrak p^n)$ 上所得的嵌入（$n\geq0$）。
对 $\mathfrak g=\mathfrak h\mathfrak p^n$，考虑 1.2.3 中的映射 $(1)_{a,\mathfrak g}$\footnote{译注：引用号 1.2.3 原书如此；该映射实际定义于 1.2.2，此处照录。}；
仍用 $\iota_n$ 表示
\[
H^0(K(\mathfrak h\mathfrak p^n)\otimes\mathbb R,\mathbb Z_p)
=\bigoplus_\iota H^0(\mathbb C,\mathbb Z_p)
\]
中的如下元素：其 $\iota$ 分量在 $\iota=\iota_n$ 时为 $1$，否则为 $0$。
于是当 $n$ 变化时，$(\iota_n)_n$ 构成一个逆系统。它在
\[
\begin{gathered}
H^0(K(\mathfrak g)\otimes\mathbb R,\mathbb Z_p)
\xrightarrow{(1)_{a,\mathfrak g}}
H^1(\mathcal O_{K(\mathfrak g)}[1/S],\mathbb Z_p(1))\\
\longrightarrow H^1((K'_{\mathfrak p})_n,\mathbb Z_p(1)),
\qquad\mathfrak g=\mathfrak h\mathfrak p^n\ (n\geq1),
\end{gathered}
\]
下的像，等于\term{范相容系统}{norm compatible system}{système compatible aux normes}
$(u_n)_n$，其中
\[
u_n=\iota_n^{-1}\bigl(\theta_a(\exp(\pi^{-n}h^{-1}\delta))\bigr),
\]
而 $h$ 是事先固定的理想 $\mathfrak h$ 的一个生成元。
令 $T=T_pE$。\translatornote{17}利用同构
\[
H_1(E(\mathbb C),\mathbb Z)\otimes\mathbb Z_p
\simeq H^0(\mathbb C,T)\simeq H^0(K^{\mathrm{ab},S},T),
\]
将 $h^{-1}\delta$ 的像定义为 $H^0(K^{\mathrm{ab},S},T)$ 的一组
$\mathcal O_{\mathfrak p}$ 基 $\xi$；最后一个同构由上面固定的
$K^{\mathrm{ab},S}\hookrightarrow\mathbb C$ 诱导。
于是，与 $u=(u_n)_n$ 和 $\xi$ 相伴的\term{Coleman幂级数}{Coleman power series}{série de Coleman}
$g_{u,\xi}$ 是 Lubin--Tate 群上的函数
\[
z\longmapsto\theta_a\bigl(z+\iota^{-1}(\exp(h^{-1}\delta))\bigr).
\]
\origpage{163}{158}
现在证明 1.2.6。在 1.2.5 中，令 $L=K(\mathfrak h\mathfrak p^n)$，$n\geq1$。
只需在 $A=\mathbb Q[\Gal(L/K)]$\translatornote{18} 且
$\mathscr F=f_*f^*(T^{\otimes(-r)}(1))$（某个 $r\geq1$）的情形下证明 1.2.6，
其中 $f$ 为典范态射 $\Spec(\mathcal O_L[1/S])\to\Spec(\mathcal O_K[1/S])$。
这是因为任意 $(\Lambda,\mathscr F)$ 都可由这种特殊情形的 $(A,\mathscr F)$，
经由 $A$ 上某个环 $?$，写成 $(?,\mathscr F\otimes_A?)$；这里的 $?$ 是
$K$ 的有限扩张的有限直积。在此情形中，有
\[
L_{A,S}(X,\mathscr F^*(1),s)
=\sum_{\mathfrak a}\nu(\mathfrak a)^{-r}N(\mathfrak a)^{-s}(\mathfrak a,L/K)^{-1},
\]
它是取值于 $\mathbb C[\Gal(L/K)]$ 的函数，其中 $\mathfrak a$ 遍历
$\mathcal O_K$ 的所有与 $S$ 互素的非零理想，$(\mathfrak a,L/K)\in\Gal(L/K)$ 表示 Artin 符号。
考虑复合映射
\[
\begin{gathered}
\Sigma\longrightarrow H^1(\mathcal O_K[1/S],\mathscr F)
\simeq H^1(\mathcal O_L[1/S],T^{\otimes(-r)}(1))\\
\longrightarrow H^1((K'_{\mathfrak p})_n,T^{\otimes(-r)}(1)),
\end{gathered}
\]
其中第一个箭头由 1.2.4 中的元素 $z_{a,\Lambda}(X,\mathscr F)$ 定义。
考虑 $\Sigma$ 中的元素，其 $\iota$ 分量在 $\iota=\iota_n$ 时为
$\delta^{\otimes(-r)}$，否则为 $0$。这个元素经上述复合映射映为
\[
h^{-r}s_r(u,\xi)\in H^1((K'_{\mathfrak p})_n,T^{\otimes(-r)}(1)).
\]
注意，我们已经算出了 Coleman 幂级数 $g_{u,\xi}$。
结合 1.1.6 中介绍的定理及第二章 2.1.7，便得到 1.2.6。

以几段略带趣味的话结束本节。

\noindent\textbf{1.2.7.} Zeta 值定义在 $\mathbb C$ 中，而显式互反律深居 $p$ 进世界的内部，
两者相隔甚远。它们之间的联系令我如此惊讶，以至于我忍不住想：这种联系是在宇宙诞生之前形成的，
\origpage{164}{159}
还是在那之后形成的？

\noindent\textbf{1.2.8.} 正如本节及第一章 §3 中所见，分圆单位和椭圆单位以许多种方式
（多得令人感到神秘！）与 zeta 值相联系。这令我惊讶，尤其是它们竟能借助显式互反律表达 zeta 值。
在 $p$ 进世界中，要表达复 zeta 函数的值，是件极其困难的事。然而，这样的困难似乎难不倒这些伟大的单位，
它们仿佛怀着强烈的愿望，想把 zeta 值表达出来。身处 $p$ 进世界，它们似乎依然记得，自己在阿基米德世界中
曾经如此伟大，曾经表达过 zeta 值。

因此，我觉得分圆单位和椭圆单位是 zeta 值的“化身”。在 [Ki] 中，仙鹤化作少女，走进与兵的家，
向他表达谢意（参见参考文献中对 [Ki] 故事的梗概）。同样地，zeta 值化身为分圆单位和椭圆单位，
走进了显式互反律。

我敬重仙鹤这份真诚的心愿。

\noindent\textbf{1.2.9.} Zeta 值的神秘性质似乎在轻声告诉我们，宇宙也具有相同的性质：
仅凭实数，无法解释宇宙。它有 $p$ 进的性质（正如某些物理学家所主张的），
并且与某些深邃的对象相联系；为方便起见，我们将这些对象称为仙鹤、银河列车，以及 zeta 值的故乡。
我们自己也可能具有相同的性质。

Zeta 元素是否也有物理意义？

\origpage{165}{160}
\section*{参考文献}
\addcontentsline{toc}{section}{加藤讲义参考文献}
\begin{description}[style=nextline,leftmargin=3.7em,labelwidth=3.2em,font=\normalfont\small,itemsep=.6em]
\item[{[Be1]}] Beilinson, A., Higher regulators and values of $L$-functions, J. Soviet Math. 30 (1985), 2036--2070.
\item[{[Be2]}] Beilinson, A., Polylogarithm and cyclotomic elements, 预印本。
\item[{[Bl]}] Bloch, S., Lectures on algebraic cycles, Duke Univ. Math. Series (1980).
\item[{[BK]}] Bloch, S. and Kato, K., $L$-functions and Tamagawa numbers motives, in The Grothendieck Festschrift, Vol. I (1980), 334--400.
\item[{[Bo]}] Borel, A., Stable real cohomology of arithmetic groups, Ann. Sci. École Norm. Sup. 7 (1974), 235--272.
\item[{[CP]}] Coates, J. and Perrin-Riou, B., On $p$-adic $L$-functions attached to motives over $\mathbb Q$, in Advanced Studies in Pure Math. 17 (1989), 23--54.
\item[{[CW]}] Coates, J. and Wiles, A., On the conjecture of Birch and Swinnerton-Dyer, Invent. Math. 39 (1977), 223--251.
\item[{[Co]}] Coleman, R., Division values in local fields, Inv. Math. 53 (1979), 91--116.
\item[{[Da]}] Damerell, R. M., $L$-functions of elliptic curves with complex multiplication, I, Acta Arith. 17 (1970), 287--301; II, 同刊 19 (1971), 311--317.
\item[{[De1]}] Deligne, P., Théorie de Hodge II, Publ. Math. IHES 40 (1972), 5--57.
\item[{[De2]}] Deligne, P., Théorème de finitude en cohomologie $\ell$-adique, in Lecture Notes in Math. 569 (SGA $4\frac12$), Springer (1977), 233--261.
\item[{[De3]}] Valeurs de fonctions $L$ et périodes d'intégrales, Proc. Symp. Pure Math., vol. 33, Part 2, AMS (1979), 313--349.
\item[{[De4]}] Deligne, P., La conjecture de Weil II, Publ. Math. IHES 52 (1981).
\item[{[De5]}] Deligne, P., Le groupe fondamental de la droite projective moins trois points, in Galois groups over $\mathbb Q$, Springer (1989), 79--298.
\item[{[dS1]}] de Shalit, E., The explicit reciprocity law in local class field theory, Duke Math. J. (1986), 163--176.
\origpage{166}{161}
\item[{[dS2]}] de Shalit, E., Iwasawa theory of elliptic curves with complex multiplication, Academic Press (1987).
\item[{[dS3]}] de Shalit, E., The explicit reciprocity law of Bloch and Kato, 预印本。
\item[{[Fa]}] Faltings, G., Crystalline cohomology and $p$-adic Galois representations, in Algebraic Analysis, Geometry, and Number Theory, Johns Hopkins Univ. (1989), 25--80.
\item[{[Fo1]}] Fontaine, J.-M., Sur certains types de représentations $p$-adiques du groupe de Galois d'un corps local: construction d'un anneau de Barsotti-Tate, Ann. of Math. 115 (1982), 547--608.
\item[{[Fo2]}] Fontaine, J.-M., Formes différentielles et modules de Tate des variétés abéliennes sur les corps locaux, Invent. Math. 65 (1982), 379--409.
\item[{[Fo3]}] Fontaine, J.-M., Cohomologie cristalline et représentations $p$-adiques, in Lecture Notes in Math. 1016, Springer (1983), 86--108.
\item[{[FM]}] Fontaine, J.-M. and Messing, W., $p$-adic periods and $p$-adic étale cohomology, Contemporary Math. 67 (1987), 179--207.
\item[{[FP1]}] Fontaine, J.-M. and Perrin-Riou, B., Autour des conjectures de Bloch et Kato, I, C. R. Acad. Sci. Paris, t. 313, Série I (1991), 189--196; II, 同刊, 349--356; III, 同刊, 421--428.
\item[{[FP2]}] Fontaine, J.-M. and Perrin-Riou, B., Autour des conjectures de Bloch et Kato, cohomologie galoisienne et valeurs de fonctions $L$, 预印本。
\item[{[Gr]}] Grothendieck, A., Formule de Lefschetz et rationalité des fonctions $L$, in Sém. Bourbaki, vol. 1965/66, Benjamin (1966), exposé 306.
\item[{[Il]}] Illusie, L., Cohomologie de de Rham et cohomologie étale $p$-adique (Sém. Bourbaki, exposé 726), in Astérisque (1990), 325--374.
\item[{[Ja]}] Jannsen, U., On the $\ell$-adic cohomology of varieties over number fields and its Galois cohomology, in Galois groups over $\mathbb Q$, Springer (1989), 315--360.
\item[{[Ka1]}] Kato, K., The explicit reciprocity law and the cohomology of Fontaine-Messing, Bull. Soc. Math. France 119 (1991), 397--441.
\origpage{167}{162}
\item[{[Ka2]}] Kato, K., Iwasawa theory and $p$-adic Hodge theory, 预印本。
\item[{[Ka3]}] Kato, K., 写作中。
\item[{[Ki]}] Kinoshita, J., The twilight-crane (1949)（《夕鹤》，一部取材于日本传说的戏剧）。

贫穷的青年与兵救下了一只被人射伤的鹤。后来的一天夜里，一位少女来到他家。
她与他共同生活，织出美丽的布匹，却请求他在她织布时不要偷看。他们结了婚，生活得很幸福。
然而，与兵后来只关心卖掉她所织布匹赚来的钱。终于有一天夜里，他偷看她织布，
发现是一只鹤正在用自己的羽毛织布。两人幸福的日子就此结束了，因为她不得不回到天空。
\item[{[KM]}] Knudsen, F. and Mumford, D., The projectivity of the moduli space of stable curves I, Math. Scand. 39, 1 (1976), 19--55.
\item[{[Ko]}] Kolyvagin, V. A., Euler systems, The Grothendieck Festschrift, vol. 2, Birkhäuser (1990), 435--483.
\item[{[Ma]}] Mazur, B., Notes on the étale cohomology of number fields, Ann. Sci. Ec. Norm. Sup. 6 (1973), 521--556.
\item[{[MW]}] Mazur, B. and Wiles, A., Class fields of abelian extensions of $\mathbb Q$, Invent. Math. 76 (1984), 179--330.
\item[{[Mi]}] Miyazawa, K.（日本诗人宫泽贤治），A night on the galaxy train（《银河铁道之夜》，约写于1924年）。

一个男孩在梦中与朋友一起乘火车游历夜空。他们沿着银河，探访北十字、天蝎、南十字等星座。
故事结尾才揭示，这位朋友已经死去；故事也暗示，这列火车正在把他送往天堂。

从数学的角度来看，男孩尚未抵达 zeta 值的故乡就从梦中醒来，但他从旅途中
（也就是从对 zeta 值的研究中）学到了许多重要的东西。
\item[{[Qu]}] Quillen, D., Higher algebraic $K$-theory, I, in Lecture Notes in Math. 341, Springer (1973), 85--147.
\item[{[Ra]}] Rapoport, M., Schappacher, N. and Schneider, P. (eds.), Beilinson's conjectures on special values of $L$-functions, Academic Press (1988).
\origpage{168}{163}
\item[{[Ru]}] Rubin, K., The ``main conjectures'' of Iwasawa theory for imaginary quadratic fields, Invent. math. 103 (1991), 25--68.
\item[{[Se]}] Serre, J.-P., Cohomologie Galoisienne, Lecture Notes in Math. 5, Springer (1965).
\item[{[So1]}] Soulé, C., $K$-théorie des anneaux d'entiers de corps de nombres et cohomologie étale, Invent. Math. 55 (1979), 251--295.
\item[{[So2]}] Soulé, C., On higher $p$-adic regulators, in Lecture Notes in Math. 854, Springer (1981), 371--401.
\item[{[So3]}] Soulé, C., The rank of étale cohomology of varieties over $p$-adic or number fields, Comp. Math. 53 (1984), 113--131.
\item[{[So4]}] Soulé, C., $p$-adic $K$-theory of elliptic curves, Duke Math. J. 54 (1987), 249--269.
\item[{[Ta1]}] Tate, J., On the conjecture of Birch and Swinnerton-Dyer and a geometric analog, in Sém. Bourbaki, vol. 1965/66, Benjamin (1966), exposé 306.
\item[{[Ta2]}] Tate, J., $p$-divisible groups, Proceedings of a conference on local fields, Driebergen, 1966, Springer (1967), 158--183.
\item[{[Wa]}] Washington, L. C., Introduction to cyclotomic fields, Springer (1982).
\item[{[We1]}] Weil, A., Elliptic functions according to Eisenstein and Kronecker, Springer (1976).
\item[{[We2]}] Weil, A., Number theory: An approach through history; From Hammurapi to Legendre, Birkhäuser (1983).
\item[{[Wi1]}] Wiles, A., Higher explicit reciprocity laws, Ann. Math. 107 (1978), 235--254.
\item[{[Wi2]}] Wiles, A., The Iwasawa conjecture for totally real fields, Ann. of Math. 131 (1990), 493--540.
\item[{[Wo]}] Wolfgang, K. S., $\lambda$-rings and Adams operators in algebraic $K$-theory, 收入 [Ra], 93--102.
\item[{[SGA4]}] Artin, M., Grothendieck, A. and Verdier, J. L., Théorie des topos et cohomologie étale des schémas, Lecture Notes in Math. 269, 270, 305, Springer (1972/73).
\end{description}
\begin{flushright}东京工业大学\end{flushright}
